\section*{第 \ref{ch:b3:conformal} 章 --- 共形映射与
黎曼映射定理}

\begin{solution}{exo:b3:conformal:1}
(a) $\varphi$ 与 $\psi(w) = \iu\frac{1+w}{1-w}$ 复合后
在两种次序下都是恒等映射（直接计算）；$\varphi$
将 $\mathbb H$ 映入 $\mathbb D$，$\psi$ 则反向映回
（\cref{ex:b3:conformal:mobius}）。取值：$\varphi(\iu) = 0$，
$\varphi(0) = -1$，$\varphi(1) = \frac{1 - \iu}{1 + \iu} =
-\iu$，且 $\varphi(z) \to 1$，当 $z \to \infty$ 时。

(b) 圆与直线是下列方程的解集
$\alpha\abs z^2 + \bar\beta z + \beta\bar z + \gamma = 0$
（$\alpha, \gamma \in \R$，$\beta \in \C$，$\abs\beta^2 >
\alpha\gamma$）：$\alpha \neq 0$ 时为圆，$\alpha = 0$ 时为直线。
代入 $z = 1/w$，并乘以 $\abs w^2$：
$\gamma\abs w^2 + \beta w + \bar\beta\bar w + \alpha = 0$ ---
--- 同一族。仿射映射显然保持该族，且
每个莫比乌斯映射都是仿射映射与一次
反演的复合（$\frac{az+b}{cz+d} = \frac ac + \frac{bc -
ad}{c}\cdot\frac1{cz + d}$，当 $c \neq 0$ 时）。
\end{solution}

\begin{solution}{exo:b3:conformal:2}
(a) 由施瓦茨引理得 $\abs{f(a)} \leq \abs a$，且等号成立（两边
均为 $= \abs a$）：取等情形迫使 $f(z) =
\eu^{\iu\theta}z$，再由 $f(a) = a$ 得 $\eu^{\iu\theta} = 1$。

(b) 设 $a \neq b$ 为不动点，并令 $g =
\varphi_a\circ f\circ\varphi_{-a} \in
\operatorname{Aut}(\mathbb D)$ --- 利用
\cref{thm:b3:conformal:autdisc}说明 $\varphi_{\pm a}$ 为自同构。则
$g(0) = \varphi_a(f(a)) = 0$，且 $g(c) = c$，其中 $c =
\varphi_a(b) \neq 0$。由 (a)，$g = \mathrm{id}$，于是 $f =
\varphi_{-a}\circ\varphi_a = \mathrm{id}$。
\end{solution}

\begin{solution}{exo:b3:conformal:3}
固定 $z$，令 $g = \varphi_{f(z)}\circ f\circ\varphi_{-z}$：
全纯映射 $\mathbb D \to \mathbb D$，且 $g(0) = 0$，于是
$\abs{g'(0)} \leq 1$（由施瓦茨引理）。利用链式法则，其中
$\varphi_a'(\zeta) = \frac{1 - \abs a^2}{(1 - \bar
a\zeta)^2}$：
\[
g'(0) = \varphi_{f(z)}'\bigl(f(z)\bigr)\cdot f'(z)\cdot
\varphi_{-z}'(0)
= \frac{1}{1 - \abs{f(z)}^2}\cdot f'(z)\cdot(1 - \abs z^2),
\]
由此得施瓦茨--皮克不等式。在某点 $z$ 处取等号时，
使 $g$ 为旋转，从而 $f =
\varphi_{-f(z)}\circ(\text{旋转})\circ\varphi_z \in
\operatorname{Aut}(\mathbb D)$ --- 于是
处处取等号（直接计算，或交换 $f, f^{-1}$
的角色后重新应用）。圆盘上的全纯自映射均是
$1$-利普希茨的，关于双曲度量 $\frac{2\abs{\dd
z}}{1 - \abs z^2}$；自同构即为其等距变换。
\end{solution}

\begin{solution}{exo:b3:conformal:4}
(a) $z \mapsto \eu^z$：将 $\{0 < \operatorname{Im}z <
\pi\}$ 双射地映到 $\mathbb H$（$\eu^{x+\iu y} =
\eu^x\eu^{\iu y}$：模不受限，辐角 $y \in
\intoo0\pi$），全纯，导数处处非零，且
逆映射也全纯（取主支 $\log$）。
(b) $z \mapsto z^2$ 使辐角加倍：开象限
$\{0 < \arg z < \frac\pi2\}$ 共形映到 $\mathbb
H$（逆映射为主支平方根）。
(c) $z \mapsto \frac{1 + z}{1 - z}$ 将 $\mathbb D$ 映到
右半平面，并保持上下对称性：它
将上半圆盘映到第一象限；
再由 (b) 平方，得到 $\mathbb H$：$z \mapsto \bigl(\frac{1
+ z}{1 - z}\bigr)^2$。
(d) 布拉施克因子 $\varphi_{1/2}(z) = \frac{z -
\frac12}{1 - \frac z2}$：$\varphi_{1/2}(\tfrac12) = 0$，且
$\varphi_{1/2}'(\tfrac12) = \frac{1 - \frac14}{(1 -
\frac14)^2} = \frac43 > 0$。
\end{solution}

\begin{solution}{exo:b3:conformal:5}
(a) 共形映射 $\C \to \mathbb D$（或 $\mathbb H$，
与凯莱映射复合后）是有界整函数：
由刘维尔定理知为常数 --- 因而不是双射。共形映射 $\mathbb
D \to \C$ 会有共形逆 $\C \to \mathbb D$：
同样矛盾。

(b) 正方形是凸的，从而单连通，且为真子集：
共形等价于 $\mathbb D$（\cref{thm:b3:conformal:rmt}）。沿负实轴割开的
平面 $\C\setminus\intoc{-\infty}0$ 关于 $1$ 为星形，
（从 $1$ 出发的线段不与割线相交），从而单连通，
且为真子集：共形等价于 $\mathbb D$。去心圆盘：若
$g \colon \mathbb D\setminus\{0\} \to \mathbb D$ 为
共形映射，则 $g$ 有界，从而 $0$ 为可去奇点
（见\cref{thm:b3:residues:riemanncw}）：$g$ 延拓为 $\tilde g
\colon \mathbb D \to \mathbb D$，且 $\tilde g(0)$ 属于
开像 $g(\mathbb D\setminus\{0\}) = \mathbb D$ 之中，
亦即某 $g(w)$，其中 $w \neq 0$；两个不相交的邻域
$0$ 与 $w$ 的邻域，其像为开集，并共享
$\tilde g(0)$，从而也共享\emph{其他}值（开集
）：$g$ 在
$\mathbb D\setminus\{0\}$ 上取某值两次 --- 与单射性矛盾。
圆环 $A = \{1 < \abs z < 2\}$：设 $F \colon \mathbb D
\to A$ 为共形映射。$F$ 在单连通的
$\mathbb D$ 上无零点，故 $F = \eu^L$，其中 $L$ 全纯
（见\cref{def:b3:conformal:simplyconnected}）。设 $\sigma$ 为
圆周 $\abs z = \frac32$ 位于 $A$ 中，并令 $\gamma =
F^{-1}\circ\sigma$，它是 $\mathbb D$ 中的闭路径；于是
\[
1 = \operatorname{Ind}_\sigma(0)
= \frac1{2\iu\pi}\int_{F\circ\gamma}\frac{\dd w}{w}
= \frac1{2\iu\pi}\int_\gamma\frac{F'}{F}
= \frac1{2\iu\pi}\int_\gamma L' = 0
\]
（$L'$ 有原函数），矛盾。去心圆盘与
圆环都不是换了装的圆盘。
\end{solution}

\begin{solution}{exo:b3:conformal:6}
(a) $T(w) = \frac{w - 1}{w + 1}$ 将 $\{\operatorname{Re}w >
0\}$ 共形地映为 $\mathbb D$（经旋转的凯莱映射：$\abs{w -
1} < \abs{w + 1}$ 当且仅当 $\operatorname{Re}w > 0$），且 $T(1) =
0$。对 $f \in \mathcal F$，$g = T\circ f\colon \mathbb D \to
\mathbb D$ 全纯，且 $g(0) = 0$；由施瓦茨引理可得
$\abs{g(z)} \leq \abs z$。

(b) 反解 $T$ 得 $f = \frac{1 + g}{1 - g}$，于是
\[
\abs{f(z)} \leq \frac{1 + \abs{g(z)}}{1 - \abs{g(z)}}
\leq \frac{1 + \abs z}{1 - \abs z} :
\]
局部一致有界（在 $\abs z \leq r < 1$ 上）。等号
在 $z_0 \neq 0$ 处成立，迫使 $\abs{g(z_0)} = \abs{z_0}$，并且
辐角一致：$g$ 为旋转，即 $f(z) = \frac{1 +
\eu^{\iu\theta}z}{1 - \eu^{\iu\theta}z}$ --- 这些是
\emph{赫格洛茨极值函数}，即共形映为右
半平面的映射。
\end{solution}

\begin{solution}{exo:b3:conformal:7}
(i) 单连通性恰好在两处用到，都是借助
无零点函数的全纯平方根的存在性
（见\cref{def:b3:conformal:simplyconnected}）：在第 0 步（
$z - b$ 的根）与第 2 步（$\varphi_a\circ f$ 的根）。
(ii) $\Omega \neq \C$ 提供了第 0 步中的点 $b$ ---
若无此条件，族 $\mathcal F$ 中将不含单射
有界映射（由刘维尔定理）。(iii) 凡恒等定理或
赫尔维茨定理（见\cref{exo:b3:residues:8}）适用之处，都需用到连通性：
所得极限为“单射或常值”，而
常值性被 $s > 0$ 排除；此外在“导数为零
则为常值”中亦用到。对 $\C$ 则失效：第 0 步不可行，且
结论亦不成立（见\cref{exo:b3:conformal:5}(a)）。
对圆环失效：圆环不单连通 --- 平方根构造
失效（例如 $z$ 本身，它在
$A$ 上无零点，却无全纯平方根，计算方式同
\cref{exo:b3:conformal:5}(b)，其中
$\frac12\int_\sigma\frac{\dd z}z \notin 2\iu\pi\Z$）--- 而
结论同样不成立。
\end{solution}

\begin{solution}{exo:b3:conformal:8}
将 $g$ 的傅里叶展开代入泊松积分，
并利用
$\frac1{2\pi}\int P_r(\varphi - t)\eu^{\iu nt}\dd t =
r^{\abs n}\eu^{\iu n\varphi}$（由 $P_r$ 的级数读出）：
$u(r\eu^{\iu\varphi}) = \sum_nc_n(g)\,r^{\abs
n}\eu^{\iu n\varphi}$，交换次序由正规收敛
保证（$\abs{c_n} \leq \norm g_\infty$，$r < 1$）。

(a) $g = \cos t$：$c_{\pm1} = \frac12$，于是 $u =
r\cos\varphi = \operatorname{Re}z = x$ --- 确为调和函数，
且边值正确。

(b) $\cos^2t = \frac12 + \frac{\cos 2t}2$：$u = \frac12 +
\frac{r^2\cos2\varphi}2 = \frac12 +
\frac{\operatorname{Re}(z^2)}2 = \frac12 + \frac{x^2 -
y^2}{2}$。

(c) $g = \mathbf 1_{(0,\pi)}$（上半圆周）：$c_0 =
\frac12$，且 $c_n = \frac{1 - (-1)^n}{2\iu\pi n}$，当 $n \neq
0$ 时，于是
\[
u(r\eu^{\iu\varphi}) = \frac12 + \frac2\pi\sum_{k\geq0}
\frac{r^{2k+1}\sin\bigl((2k+1)\varphi\bigr)}{2k + 1},
\qquad u(0) = \frac12 :
\]
圆心处的值恰为边界数据的平均 --- 平均值性质的
直接体现。
\end{solution}

\begin{solution}{exo:b3:conformal:9}
由 $(1-r)^2 \leq 1 - 2r\cos\theta + r^2 \leq (1+r)^2$ 可得：
\[
\frac{1-r}{1+r} = \frac{1 - r^2}{(1+r)^2} \leq P_r(\theta)
\leq \frac{1 - r^2}{(1 - r)^2} = \frac{1+r}{1-r} .
\]
设 $u \geq 0$ 为 $\mathbb D$ 上的调和函数，$s < 1$，记 $u_s(z) =
u(sz)$ 在 $\bar{\mathbb D}$ 的某邻域上调和，
因而等于其泊松积分
（见\cref{thm:b3:conformal:poisson}，唯一性用于其
自身边界值）；夹逼核函数，并利用
平均值性质 $\frac1{2\pi}\int u_s(\eu^{\iu t})\dd t = u(0)$：
\[
\frac{1-r}{1+r}\,u(0) \leq u(s\,r\eu^{\iu\varphi}) \leq
\frac{1+r}{1-r}\,u(0).
\]
令 $s \to 1^-$，固定 $r\eu^{\iu\varphi}$（利用
$u$ 的连续性）：得哈纳克不等式。若 $u$ 在 $\C$ 上调和，
且 $u \geq m$：对 $u - m$ 在圆盘 $D(0, R)$ 上应用哈纳克不等式，
即对 $z \mapsto u(Rz) - m$：固定 $z$，并取 $r =
\abs z/R \to 0$，两个界都趋于 $u(0) - m$：$u(z) =
u(0)$ --- 为常值（由单侧界推出的双侧
刘维尔型结论）。
\end{solution}

\begin{solution}{exo:b3:conformal:10}
局部地，$u = \operatorname{Re}F$，其中 $F$ 全纯
（见\cref{prop:b3:conformal:harmonicholo}），于是 $u\circ f =
\operatorname{Re}(F\circ f)$ 在有定义之处皆调和：
调和性是共形不变的。在 $\mathbb H$ 上，
主值对数给出 $\log z = \ln\abs z + \iu\arg z$，
它在 $\mathbb H$ 上全纯，于是 $u = \frac1\pi\arg z =
\operatorname{Im}\bigl(\frac1\pi\log z\bigr)$ 是调和的，
带有边界极限：对 $x > 0$，$\arg \to 0$，$u \to 0$；
对 $x < 0$，$\arg \to \pi$，$u \to 1$：边值为 $\mathbf
1_{\intoo{-\infty}0}$ 在每个 $x \neq 0$ 处。
经凯莱映射搬运（它在取倒数后把 $\mathbb D \to \mathbb H$，
并将上半圆周对应到负
实轴，差一个由三个边界点定出的旋转
），$u\circ(\text{凯莱})$ 解决了
\cref{exo:b3:conformal:8}(c) 中的圆盘问题；在圆心处
求值得 $u = \frac12$，而闭形式
$\frac1\pi\arg$ 可通过对级数求和来检验，
$\sum\frac{r^{2k+1}\sin((2k+1)\varphi)}{2k+1} =
\frac12\arctan\frac{2r\sin\varphi}{1 - r^2}$ 型恒等式
--- 这是得到同一答案的初等途径。
\end{solution}

\begin{solution}{exo:b3:conformal:11}
(a) 设 $a \neq b$ 固定。以 $\varphi_a(z) =
\frac{z - a}{1 - \bar az}$ 作共轭（交换 $a$
与 $0$ 的自同构），$g = \varphi_a\circ f\circ\varphi_a^{-1}$ 固定
$0$ 及点 $c = \varphi_a(b) \neq 0$。由施瓦茨引理：
$\abs{g(z)} \leq \abs z$，且在 $z = c$ 处取等号
（$g(c) = c$）：取等情形迫使 $g(z) = \lambda z$，
其中 $\abs\lambda = 1$，由 $\lambda c = c$ 得 $\lambda =
1$，即 $g = \mathrm{id}$，从而 $f = \mathrm{id}$。

(b) 令 $h(z) = f(z)/z$（$0$ 是可去奇点），
则它在 $\mathbb D$ 上全纯，且 $\abs h \leq 1$（由施瓦茨引理）；
$\abs h < 1$ 处处成立，否则由最大模原理
（$\abs h$ 在内点取最大值）得 $h$ 为模为 1 的常数，
即 $f$ 为旋转 --- 被排除。在紧集
$\bar D(0,r)$ 上 $c_r = \max\abs h < 1$，故 $\abs{f(z)} \leq
c_r\abs z$ 在其上成立；又 $f$ 把 $\bar D(0,r)$ 映入
自身（$c_r\abs z \leq r$），于是该估计可迭代：
$\abs{f^{\circ n}(z)} \leq c_r^n\,r \to 0$ 在
$\bar D(0, r)$ 上一致成立。

(c) $\frac{z^2 + z}2$ 的不动点：$z^2 + z = 2z$
当且仅当 $z(z - 1) = 0$；仅 $z = 0$ 位于 $\mathbb D$ 内（$z = 1$ 在
边界上）。不是旋转（$f'(0) = \frac12$），故
轨道趋于 $0$；定量地，$f(z) = \frac z2(1 + z)$
给出 $\abs{f(z)} \leq \frac{3}{4}\abs z$，在 $\abs z \leq
\frac12$ 上成立；一旦轨道足够小，$\abs{f(z)} \approx
\frac{\abs z}2$，渐近地为公比
$f'(0) = \frac12$ 的几何收缩。从 $z_0 = \frac12$ 出发：$z_1 = \frac38$，
$z_2 \approx 0.258$，$z_3 \approx 0.162$ --- 每步减半，
与乘子所预言的一致。
\end{solution}

\begin{solution}{exo:b3:conformal:12}
(a) $\Delta u = 6x - 6x + 0 = 0$。柯西--黎曼方程要求
$v_y = u_x = 3x^2 - 3y^2$ 与 $v_x = -u_y = 6xy - 2$。
对第一式关于 $y$ 积分：$v = 3x^2y - y^3 + c(x)$；
代入第二式：$6xy + c'(x) = 6xy - 2$，故 $c(x)
= -2x + C$。于是 $v = 3x^2y - y^3 - 2x + C$，从而
\[
f = u + \iu v = (x^3 - 3xy^2) + \iu(3x^2y - y^3)
+ 2y - 2\iu x + \iu C = z^3 - 2\iu z + \iu C .
\]

(b) 形式 $\omega = -u_y\,\dd x + u_x\,\dd y$ 是闭的，
恰因 $\Delta u = 0$（$\partial_y(-u_y) =
-u_{yy} = u_{xx} = \partial_x(u_x)$）。在星形开集上，
由庞加莱引理（\cref{thm:b3:forms:poincare}；或
由原函数的构造，
\cref{thm:b3:holomorphic:cauchy} 施于全纯函数
$u_x - \iu u_y$）给出 $v$，满足 $\dd v = \omega$，即
柯西--黎曼方程组：$u + \iu v$ 全纯。两个共轭
之差是梯度为零的函数，
即常数（由连通性）。

(c) 对 $u = \ln\abs z$（定义在 $\C^*$ 上），$\omega = -u_y\dd x +
u_x\dd y = \frac{-y\,\dd x + x\,\dd y}{x^2 + y^2} =
\omega_\theta$，即角形式
（\cref{ex:b3:forms:angular}），其沿单位
圆周的积分为 $2\pi \neq 0$：不恰当，故不存在
整体共轭 --- 共轭即辐角的一个连续
选取，而环绕数恰好就是
障碍。局部地（在任一星形子区域上），$v =
\arg z$ 成立，且 $u + \iu v = \log z$：失败
只在整体上，而非局部。
\end{solution}

\begin{solution}{pb:b3:conformal:1}
\textbf{1.} $g$ 单射且全纯；在 $C_\rho$
（$\rho > 1$）上它光滑，所围面积为
$\frac1{2\iu}\oint_{g(C_\rho)}\bar\zeta\,\dd\zeta$（本科二年级
的格林--黎曼面积公式，取正
定向）。作替换 $\zeta = g(w)$，$w =
\rho\eu^{\iu\theta}$：
\[
A_\rho = \frac1{2\iu}\int_{C_\rho}\overline{g(w)}\,g'(w)\dd
w = \frac\rho2\int_0^{2\pi}\overline{g(\rho\eu^{\iu\theta})}
\,g'(\rho\eu^{\iu\theta})\,\eu^{\iu\theta}\,\dd\theta .
\]
代入 $\bar g = \rho\eu^{-\iu\theta} + \bar b_0 +
\sum_m\bar b_m\rho^{-m}\eu^{\iu m\theta}$ 及 $g' = 1 -
\sum_nnb_n\rho^{-n-1}\eu^{-\iu(n+1)\theta}$。
乘以 $\eu^{\iu\theta}$ 后，只有零频率的
乘积在 $\theta$ 积分中保留下来（正规
收敛允许逐项计算）：这一对
$(\rho\eu^{-\iu\theta})\cdot1$ 贡献 $2\pi\rho$，而
每一对 $\bar b_n\rho^{-n}\eu^{\iu
n\theta}\cdot(-nb_n\rho^{-n-1}\eu^{-\iu(n+1)\theta})$
贡献 $-2\pi n\abs{b_n}^2\rho^{-2n-1}$：
\[
A_\rho = \pi\rho^2 - \pi\sum_{n\geq1}n\abs{b_n}^2\rho^{-2n}.
\]

\textbf{2.} 面积非负：$\sum_{n\leq
N}n\abs{b_n}^2\rho^{-2n} \leq \rho^2$ 对每个 $N$ 成立；令
$\rho \to 1^+$，再令 $N \to \infty$：
$\sum_{n\geq1}n\abs{b_n}^2 \leq 1$。在 $\abs{b_1}
\leq 1$ 取等号时迫使其余 $b_n = 0$：$g(w) = w + b_0 +
\eu^{\iu\alpha}/w$，它映到一条
长为 $4$ 的线段的补集上（茹科夫斯基型映射）：即极值函数。

\textbf{3.} $f(z)/z = 1 + a_2z + \cdots$ 全纯且
在 $\mathbb D$ 上无零点（$f$ 仅在 $0$ 处为零，且为单零
点：由单射性），故 $f(z^2)/z^2$ 亦然，它有
全纯平方根 $\varphi$，满足 $\varphi(0) = 1$
（\cref{def:b3:conformal:simplyconnected}；$\mathbb D$ 为
凸集）。于是 $h(z) = z\varphi(z^2)$ 满足 $h^2 =
f(z^2)$，$h(z) = z\bigl(1 + \frac{a_2}2z^2 + \cdots\bigr)$
（取根用二项式级数），且 $h(z) = z\varphi(z^2)$ 由构造是
奇函数（$\varphi(z^2)$ 为偶函数）。单射性：由 $h(z) = h(z')$ 得 $f(z^2) =
f(z'^2)$，故 $z^2 = z'^2$，即 $z' = \pm z$；若 $z' = -z$
则由奇性得 $h(z) = -h(z)$，从而 $h(z) = 0$，迫使 $z
= 0$（$\varphi$ 无零点）：$z = z' = 0$。

\textbf{4.} 令 $g(w) = 1/h(1/w)$：当 $\abs w > 1$ 时，$1/w \in
\mathbb D\setminus\{0\}$，且此处 $h \neq 0$：良定义，
单射（单射的复合），展开式为
\[
g(w) = \frac{1}{\frac1w\bigl(1 + \frac{a_2}{2w^2} +
\cdots\bigr)} = w\Bigl(1 - \frac{a_2}{2w^2} + \cdots\Bigr)
= w - \frac{a_2}{2}\,w^{-1} + \cdots :
\]
此为第一部分的形式，其中 $b_1 = -\frac{a_2}2$。面积定理
给出 $\abs{\frac{a_2}2} \leq 1$，即 $\abs{a_2} \leq 2$。

\textbf{5.} $k(z) = \frac z{(1-z)^2} = z\sum_{m\geq0}(m +
1)z^m = \sum_{n\geq1}nz^n$：系数为 $a_n = n$，故 $a_2 =
2$。单叶性与像：$k = \frac14\bigl[w^2 - 1\bigr]$
其中 $w = \frac{1 + z}{1 - z}$，将 $\mathbb D$
共形映为右半平面；$w^2$ 将该半平面
共形映为 $\C\setminus\intoc{-\infty}0$；再经
$\frac{\cdot - 1}4$ 得
$\C\setminus\intoc{-\infty}{-\frac14}$：每一步均单射，
像如所述。

\textbf{6.} 令 $F = \frac{cf}{c - f}$。因 $c \notin f(\mathbb
D)$，分母处处不为零，故 $F$ 全纯，并且
单射（$w \mapsto \frac{cw}{c - w}$ 为莫比乌斯变换，
在 $w = c$ 之外单射）。展开时取 $f = z + a_2z^2 + \cdots$，
\[
F = f\cdot\frac1{1 - f/c} = \bigl(z + a_2z^2\bigr)\Bigl(1 +
\frac zc\Bigr) + O(z^3)
= z + \Bigl(a_2 + \frac1c\Bigr)z^2 + O(z^3):
\]
$F \in \mathcal S$，且 $A_2 = a_2 + \frac1c$。

\textbf{7.} 两次应用比伯巴赫定理：$\abs{a_2} \leq 2$ 且
$\abs{a_2 + \frac1c} \leq 2$，故 $\abs{\frac1c} \leq 4$：
$\abs c \geq \frac14$。不被取到的每个值都位于
$D(0,\frac14)$ 之外，即 $f(\mathbb D) \supseteq D(0, \frac14)$。
此界为最优：科贝函数不取 $-\frac14$（第 5 题）。

\textbf{8.} 设 $d = d(z_0, \partial\Omega)$，并令 $\psi =
\varphi^{-1} \colon \mathbb D \to \Omega$, $\psi(0) = z_0$,
$\psi'(0) = 1/\varphi'(z_0) > 0$。归一化 $\tilde
f(z) = \frac{\psi(z) - z_0}{\psi'(0)}$ 属于 $\mathcal S$，
故其像包含 $D(0, \frac14)$；还原尺度后，
$\Omega = \psi(\mathbb D) \supseteq D\bigl(z_0,
\tfrac{\psi'(0)}4\bigr)$，故 $d \geq \frac{\psi'(0)}4 =
\frac1{4\varphi'(z_0)}$：左侧不等式。右侧：$\varphi\circ(z_0 + d\,\cdot)$
将 $\mathbb D$ 映入 $\mathbb D$，且 $0 \mapsto 0$；施瓦茨引理
给出其在 $0$ 处的导数界：$d\,\varphi'(z_0) \leq 1$，即
$\frac1{\varphi'(z_0)} \geq d$ --- 合并得
\[
d \leq \frac1{\varphi'(z_0)} \leq 4d .
\]
（题中所示左侧不等式就是上述链式
重新整理后的形式。）

\textbf{9.} 对科贝函数 $a_n = n$：
猜想中处处取等；其旋转
$\eu^{-\iu\theta}k(\eu^{\iu\theta}z) = \sum
n\eu^{\iu(n-1)\theta}z^n$ 亦有 $\abs{a_n} = n$。沿用
第 4 题的思路：$h = z + \frac{a_2}2z^3 + c_5z^5 + \cdots$，且满足
$h^2 = f(z^2)$；比较 $z^6$ 的系数：$2c_5 +
\frac{a_2^2}4 = a_3$，故 $c_5 = \frac{a_3}2 -
\frac{a_2^2}8$。于是
\[
g(w) = \frac1{h(1/w)} = w\Bigl(1 - \frac{a_2}{2w^2} +
\Bigl(\frac{a_2^2}4 - c_5\Bigr)\frac1{w^4} + \cdots\Bigr)
= w - \frac{a_2}2w^{-1} + \Bigl(\frac{3a_2^2}8 -
\frac{a_3}2\Bigr)w^{-3} + \cdots,
\]
于是由面积定理（$\sum n\abs{b_n}^2 \leq 1$）得
\[
\Bigl|\frac{a_2}2\Bigr|^2 + 3\,\Bigl|\frac{3a_2^2}8 -
\frac{a_3}2\Bigr|^2 \leq 1 .
\]
科贝函数检验（$a_2 = 2$，$a_3 = 3$）：$\abs{\frac{a_2}2}^2 = 1$
且 $\frac{3\cdot4}8 - \frac32 = 0$，总和恰为 $1$ ---
取极值，理应如此。

\textbf{10.} $\varphi$ 为圆盘自同构，且
$\varphi(0) = z_0$，故 $f\circ\varphi$ 在
$\mathbb D$ 上单叶（单射的复合），且 $f'(z_0) \neq
0$（见\cref{def:b3:conformal:conformal}）：$F$ 良定义
且单叶。$F(0) = 0$。由商的求导法则：
\[
\varphi'(z) = \frac{(1 + \bar z_0z) - (z + z_0)\bar z_0}
{(1 + \bar z_0z)^2}
= \frac{1 - \abs{z_0}^2}{(1 + \bar z_0z)^2},
\qquad \varphi'(0) = 1 - \abs{z_0}^2 ,
\]
故 $(f\circ\varphi)'(0) = f'(z_0)(1 - \abs{z_0}^2)$，且
$F'(0) = 1$，故 $F \in \mathcal S$。

\textbf{11.} 记 $G = f\circ\varphi$，则 $G'' =
f''(\varphi)\,\varphi'^2 + f'(\varphi)\,\varphi''$，且
$\varphi''(z) = -2\bar z_0(1 - \abs{z_0}^2)(1 + \bar
z_0z)^{-3}$ 给出 $\varphi''(0) = -2\bar z_0(1 -
\abs{z_0}^2)$。因此
\[
A_2 = \frac{G''(0)}{2f'(z_0)(1 - \abs{z_0}^2)}
= \frac12\Bigl[\bigl(1 - \abs{z_0}^2\bigr)
\frac{f''(z_0)}{f'(z_0)} - 2\bar z_0\Bigr] .
\]
对 $F$ 应用比伯巴赫（第 4 题）：$\abs{A_2} \leq 2$，即
$\bigl|(1 - \abs{z_0}^2)\frac{f''(z_0)}{f'(z_0)} - 2\bar
z_0\bigr| \leq 4$。两边同乘
$z_0/(1 - \abs{z_0}^2)$，其模为
$r/(1 - r^2)$，其中 $z_0 = r\eu^{\iu\theta}$，并利用 $z_0\bar
z_0 = r^2$：
\[
\Bigl|\,z_0\frac{f''(z_0)}{f'(z_0)} -
\frac{2r^2}{1 - r^2}\Bigr| \leq \frac{4r}{1 - r^2} .
\]

\textbf{12.} 若复数与某实点的距离不超过 $\rho$，
则其实部落在以 $c$ 为中心的区间 $\intcc{c - \rho}{c +
\rho}$ 中：将此结论应用于 $c = \frac{2r^2}{1-r^2}$，$\rho =
\frac{4r}{1-r^2}$。

\textbf{13.} 对无零点的 $\mathcal C^1$ 函数 $g$：
$\log\abs g = \frac12\log(g\bar g)$，故 $\frac{\dd}{\dd
t}\log\abs g = \frac{g'\bar g + g\bar g'}{2\abs g^2} =
\operatorname{Re}\frac{g'}g$。令 $g(t) =
f'(t\eu^{\iu\theta})$（无零点：由单叶性），$g'(t) =
\eu^{\iu\theta}f''(t\eu^{\iu\theta})$，于是在 $z =
t\eu^{\iu\theta}$：
\[
\frac{\dd}{\dd t}\log\bigl|f'(t\eu^{\iu\theta})\bigr|
= \operatorname{Re}\Bigl(\eu^{\iu\theta}
\frac{f''}{f'}\Bigr)
= \frac1t\operatorname{Re}\Bigl(z\frac{f''}{f'}\Bigr)
\in \Bigl[\frac{2t - 4}{1 - t^2},\
\frac{2t + 4}{1 - t^2}\Bigr]
\]
由第 12 题（取半径 $t$）可得。由于 $\frac{\dd}{\dd
t}\log\frac{1+t}{(1-t)^3} = \frac1{1+t} + \frac3{1-t} =
\frac{2t+4}{1-t^2}$，且 $\frac{\dd}{\dd
t}\log\frac{1-t}{(1+t)^3} = \frac{-1}{1-t} - \frac3{1+t} =
\frac{2t-4}{1-t^2}$，在 $0$ 到 $r$ 上积分（三个
函数在 $t = 0$ 处均为零，且 $f'(0) = 1$），得
\[
\log\frac{1-r}{(1+r)^3} \leq \log\abs{f'(z)} \leq
\log\frac{1+r}{(1-r)^3} :
\]
取指数即得畸变定理。

\textbf{14.} 上界：沿线段 $[0, z]$，
\[
\abs{f(z)} = \Bigl|\int_0^rf'(t\eu^{\iu\theta})
\eu^{\iu\theta}\dd t\Bigr|
\leq \int_0^r\frac{1+t}{(1-t)^3}\dd t
= \frac{r}{(1-r)^2}
\]
（$\frac{\dd}{\dd t}\frac t{(1-t)^2} = \frac{1+t}{(1-t)^3}$）。
下界：注意到 $\frac r{(1+r)^2} < \frac14$ 对 $r < 1$ 成立（因为
$(1-r)^2 > 0$），故若 $\abs{f(z)} \geq \frac14$，则
无需证明。否则线段 $[0, f(z)]$ 位于
$D(0, \frac14) \subseteq f(\mathbb D)$ 中（第 7 题），且
$\gamma = f^{-1}\circ[0, f(z)]$ 是 $\mathcal C^1$ 路径，
自 $0$ 到 $z$，位于 $\mathbb D$ 中（$f^{-1}$ 全纯，
见\cref{cor:b3:residues:openmapping}）。代入 $w =
f(\zeta)$，
\[
\abs{f(z)} = \int_{[0,f(z)]}\abs{\dd w}
= \int_\gamma\abs{f'(\zeta)}\,\abs{\dd\zeta}
\geq \int_0^1\frac{1 - \rho(s)}{(1 + \rho(s))^3}
\,\abs{\gamma'(s)}\,\dd s
\]
其中 $\rho = \abs\gamma$，利用
半径为 $\rho(s)$ 处的下界畸变估计。由于 $\abs{\gamma'} \geq \rho'$
（在有定义处成立；$\rho$ 为利普希茨函数），且被积函数中的
因子为正，故
\[
\abs{f(z)} \geq \int_0^1
\frac{1 - \rho(s)}{(1 + \rho(s))^3}\,\rho'(s)\,\dd s
= \int_0^{r}\frac{1 - u}{(1 + u)^3}\,\dd u
= \frac{r}{(1 + r)^2}
\]
（$\rho(0) = 0$，$\rho(1) = r$；此代换只用到
一个原函数，即 $\frac{\dd}{\dd u}\frac u{(1+u)^2} =
\frac{1-u}{(1+u)^3}$，而不需要单调性）。

\textbf{15.} $k'(z) = \frac{\dd}{\dd z}\,z(1 - z)^{-2} =
\frac{1 + z}{(1 - z)^3}$。当 $z = r$ 时：$k'(r) =
\frac{1+r}{(1-r)^3}$，$k(r) = \frac r{(1-r)^2}$ --- 两者
均取到上界。当 $z = -r$ 时：$k'(-r) =
\frac{1-r}{(1+r)^3}$，$\abs{k(-r)} = \frac r{(1+r)^2}$
--- 两者均取到下界。克贝函数将
其正实轴最大程度地拉伸至远处边界，并
将对径射线最大程度地压向其割缝的
尖端。

\textbf{16.} $\abs{a_2} = 2$ 意味着 $\abs{b_1} = 1$，对 $g(w)
= 1/h(1/w) = w - \frac{a_2}2w^{-1} + \cdots$（第 4 题）；
面积定理 $\sum n\abs{b_n}^2 \leq 1$ 于是使
其余所有系数为零：$g(w) = w + b_0 + \eu^{\iu\alpha}/w$，其中
$\eu^{\iu\alpha} = -\frac{a_2}2$。由于 $h$ 为奇函数，$g(-w) =
1/h(-1/w) = -1/h(1/w) = -g(w)$：$g$ 为奇函数，故 $b_0 = 0$。
取逆得 $h(z) = 1/g(1/z) = \frac{z}{1 +
\eu^{\iu\alpha}z^2}$，再由 $f(z^2) = h(z)^2$ 得
\[
f(w) = \frac{w}{(1 + \eu^{\iu\alpha}w)^2}
= \frac{w}{(1 - \eu^{\iu\theta}w)^2}
= \eu^{-\iu\theta}k\bigl(\eu^{\iu\theta}w\bigr),
\qquad \eu^{\iu\theta} = -\eu^{\iu\alpha} .
\]
反之，每个旋转均有 $a_2 = 2\eu^{\iu\theta}$，其
模为 $2$：比伯巴赫不等式的极值函数恰为
旋转后的克贝函数。

\textbf{17.} 若 $c$ 为被省略值，且 $\abs c = \frac14$：
由第 6 题，存在 $F \in \mathcal S$，使 $A_2 = a_2 +
\frac1c$，故 $\frac1c = A_2 - a_2$，且
\[
4 = \Bigl|\frac1c\Bigr| = \abs{A_2 - a_2}
\leq \abs{A_2} + \abs{a_2} \leq 2 + 2 = 4 :
\]
处处取等，特别地 $\abs{a_2} = 2$。由
第 16 题，$f = \eu^{-\iu\theta}k(\eu^{\iu\theta}\cdot)$，
其被省略值集合为 $\eu^{-\iu\theta}
\intoc{-\infty}{-\frac14}$：
模为 $\frac14$ 的唯一的被省略值为 $c = -\eu^{-\iu\theta}/4$。（验算：
$a_2 = 2\eu^{\iu\theta}$，$\frac1c = -4\eu^{\iu\theta}$，
故 $A_2 = -2\eu^{\iu\theta} = -a_2$：三角不等式
以反向对齐的方式取等，这是必然的。）

\textbf{18.} 圆周 $\abs z =
r$ 上的柯西估计（\cref{thm:b3:holomorphic:analytic}），结合
增长定理，可得
\[
\abs{a_n} \leq r^{-n}\sup_{\abs z = r}\abs f
\leq r^{1-n}\,(1 - r)^{-2} .
\] 取 $r = 1 - \frac1n$（$n \geq
2$）：$r^{1-n} = \bigl(1 + \frac1{n-1}\bigr)^{n-1} < \eu$
（单调递增数列，极限为 $\eu$），且 $(1 - r)^{-2} =
n^2$：$\abs{a_n} < \eu\,n^2$。

\textbf{19.} $U = f(D(0, r))$ 为开集
（见\cref{cor:b3:residues:openmapping}），且包含 $0$。
边界：若 $p \in \partial U$，则可写 $p = \lim f(z_k)$
其中 $z_k \in D(0, r)$；取子列得 $z_k \to z_\infty
\in \bar D(0, r)$，且 $f(z_\infty) = p \notin U$ 迫使
$z_\infty \in \partial D(0, r)$：$\partial U \subseteq
f(\partial D(0, r))$，故 $U$ 的每个边界点之模
均 $\geq \frac r{(1+r)^2}$（第 14 题）。现设
$\abs w < \frac r{(1+r)^2}$，且 $w \notin U$。线段
$[0, w]$ 连通，且与 $U$（在 $0$ 处）
及其补集（在 $w$ 处）相交，故它与 $\partial U$ 相交；但其上各点
的模均 $\leq \abs w < \frac r{(1+r)^2}$：
矛盾。故 $D\bigl(0, \frac r{(1+r)^2}\bigr)
\subseteq U$。锐性：$k(-r) = -\frac r{(1+r)^2}$ 为
$D(0,r)$ 边界点的像，且 $k$ 是单射，
故 $k(-r) \notin k(D(0, r))$：半径不能
再增大。当 $r \to 1^-$ 时，$\frac r{(1+r)^2} \to \frac14$：
这即整个单位圆盘上的四分之一定理。

\textbf{20.} 首先，$\partial f(\mathbb D) \neq \emptyset$：
否则 $f(\mathbb D)$ 为非空的既开又闭集，将等于
整个 $\C$，而 $f^{-1} \colon \C \to \mathbb D$ 将是一个
有界非常值整函数，这与
\cref{cor:b3:holomorphic:liouville}相矛盾。左侧不等式：
第 10 题中的科贝变换 $F$ 属于 $\mathcal S$，且
$F(\mathbb D) = \frac{f(\mathbb D) -
f(z_0)}{f'(z_0)(1-\abs{z_0}^2)}$；由第 7 题，它包含
$D(0, \frac14)$，因此
\[
f(\mathbb D) \supseteq f(z_0) + D\Bigl(0,\
\tfrac14\abs{f'(z_0)}\bigl(1 - \abs{z_0}^2\bigr)\Bigr) ,
\]
而 $\partial f(\mathbb D)$ 的每一点 ---
与开集 $f(\mathbb D)$ 不交 --- 距离 $\geq
\frac14(1-\abs{z_0}^2)\abs{f'(z_0)}$，以 $f(z_0)$ 为基点。右侧
不等式：设 $d = d(f(z_0), \partial f(\mathbb D)) <
\infty$。圆盘 $D(f(z_0), d)$ 含于 $f(\mathbb D)$：从
$f(z_0)$ 到其中任一点的线段，其上各点的距离皆
$< d$（以 $f(z_0)$ 为基点），故永不与 $\partial f(\mathbb D)$ 相交，
由第 19 题的连通性论证，它留在
$f(\mathbb D)$ 中。于是 $\hat g(w) = f^{-1}(f(z_0) + dw)$ 将
$\mathbb D$ 映入 $\mathbb D$，且 $\hat g(0) = z_0$，又
$\chi = \psi^{-1}\circ\hat g$，其中 $\psi(z) = \frac{z +
z_0}{1 + \bar z_0z}$，它固定 $0$，由施瓦茨
引理（见\cref{thm:b3:conformal:schwarz}）给出 $\abs{\chi'(0)}
\leq 1$。由于 $\chi'(0) = \frac{\hat g'(0)}{\psi'(0)} =
\frac{d}{f'(z_0)\,(1 - \abs{z_0}^2)}$，即 $d \leq (1 -
\abs{z_0}^2)\abs{f'(z_0)}$。

\textbf{21.} $k(\mathbb D) = \C\setminus
\intoc{-\infty}{-\frac14}$，故 $\partial k(\mathbb D) =
\intoc{-\infty}{-\frac14}$，且对正实点
$k(r) = \frac r{(1-r)^2}$，其最近的边界点为
$-\frac14$：
\[
d = k(r) + \frac14 = \frac{4r + (1-r)^2}{4(1-r)^2}
= \frac{(1+r)^2}{4(1-r)^2} .
\]
第 20 题的左端：$\frac14(1 - r^2)k'(r) =
\frac14\,\frac{(1-r^2)(1+r)}{(1-r)^3} =
\frac{(1+r)^2}{4(1-r)^2} = d$：等号成立。（右端
等于 $4d$：整体因子 $4$ 将两端区分开，
而科贝函数恰好取到下界。）

\textbf{22.} 唯一的原理是帕塞瓦尔恒等式：单叶性
禁止重叠，因而
$\{\abs w > \rho\}$ 的像的补集面积，在圆周上按傅里叶模展开，
非负 --- 面积定理即带符号的 $L^2$
恒等式。其余内容皆为此不等式
之迁移：开平方（第 3 题）将其化为
$\abs{a_2} \leq 2$；由省略值出发的莫比乌斯反射
（第 6 题）将 $\abs{a_2} \leq 2$ 化为
四分之一定理；圆盘自同构（第 10 题）将 $\abs{a_2}
\leq 2$ 扩展到整个圆盘，得畸变定理；
径向积分将畸变化为增长，并
将增长化为覆盖。在每一阶段，等号情形
也同样传递，总落在旋转科贝
函数 --- 全理论只有一族极值函数，
正如旋转恰为施瓦茨
引理的唯一极值。一个内部不等式，加上其取等情形的刚性，
支配整个几何：共形
映射正是对此类不等式加以利用的艺术。

\textbf{23.} 增长定理给出 $\abs f \leq
\frac{r}{(1-r)^2}$ 在每个 $\bar D(0, r)$ 上一致成立，$r <
1$，对所有 $f \in \mathcal S$ 同时成立：局部有界，故
$\mathcal S$ 为正规族
（见\cref{thm:b3:conformal:montel}）。若 $f_n \in \mathcal S \to
f$ 局部一致收敛：$f$ 全纯，且 $f_n' \to f'$
局部一致（魏尔斯特拉斯），故 $f(0) = 0$，$f'(0) = 1$
--- 特别地 $f$ 非常值 --- 由胡尔维茨
定理（见\cref{exo:b3:residues:8}）使单射之极限
仍为单射：$f \in \mathcal S$。连续泛函（如
$f \mapsto \abs{a_2} = \frac{\abs{f''(0)}}2$）在紧类上
\emph{取到}其上确界：极值函数
即已存在，远在人们知道其形式之前 ---
这正是一切针对系数问题的变分论证的
出发点，比伯巴赫的问题亦在其中。

\textbf{24.} 记 $g(w) = w + A_2w^2 + O(w^3)$，复合得：
\[
z = g(f(z)) = f(z) + A_2f(z)^2 + O(z^3)
= z + (a_2 + A_2)z^2 + O(z^3),
\]
故 $A_2 = -a_2$，且 $\abs{A_2} \leq 2$（比伯巴赫），
等号成立当且仅当 $\abs{a_2} = 2$，即 $f$ 为旋转后的
科贝函数（第 16 题）--- 此时 $g$ 是
对应的逆，定义于缝隙平面上。

\textbf{25.} $f(z_1) - f(z_2) = (z_1 - z_2)\bigl(1 + a(z_1 +
z_2)\bigr)$。若 $\abs a \leq \frac12$，则对 $z_1 \neq z_2$ 位于
$\mathbb D$ 中，$\abs{a(z_1 + z_2)} < 2\abs a \leq 1$（严格：
$\abs{z_1 + z_2} < 2$），故第二个因子不可能为零：
单射，且 $f \in \mathcal S$（归一化条件已
包含在内）。若 $\abs a > \frac12$，则点 $s =
-\frac1a$ 的模 $\abs s < 2$，于是 $z_{1,2} = \frac s2 \pm
\varepsilon$ 位于 $\mathbb D$ 中（$\varepsilon > 0$ 充分小），
两者互异，且 $z_1 + z_2 = s$ 使该因子为零：$f(z_1)
= f(z_2)$，不是单射。因此 $\mathcal S$ 含有 $z +
az^2$ 当且仅当 $\abs a \leq \frac12$。比伯巴赫的界
$\abs{a_2} \leq 2$ 因而远未被
次数为 $2$ 的多项式取到 --- 科贝函数的
系数 $a_n = n$ 来自一个无穷级数，
沿省略射线协同作用，这种性态是任何多项式
（仅在系数很小时才属于 $\mathcal S$）
所无法模仿。
\end{solution}
